% Options for packages loaded elsewhere
\PassOptionsToPackage{unicode}{hyperref}
\PassOptionsToPackage{hyphens}{url}
\documentclass[
  11pt,
]{article}
\usepackage{xcolor}
\usepackage[margin=1in]{geometry}
\usepackage{amsmath,amssymb}
\setcounter{secnumdepth}{-\maxdimen} % remove section numbering
\usepackage{iftex}
\ifPDFTeX
  \usepackage[T1]{fontenc}
  \usepackage[utf8]{inputenc}
  \usepackage{textcomp} % provide euro and other symbols
\else % if luatex or xetex
  \usepackage{unicode-math} % this also loads fontspec
  \defaultfontfeatures{Scale=MatchLowercase}
  \defaultfontfeatures[\rmfamily]{Ligatures=TeX,Scale=1}
\fi
\usepackage{lmodern}
\ifPDFTeX\else
  % xetex/luatex font selection
  \setmainfont[]{Times New Roman}
\fi
% Use upquote if available, for straight quotes in verbatim environments
\IfFileExists{upquote.sty}{\usepackage{upquote}}{}
\IfFileExists{microtype.sty}{% use microtype if available
  \usepackage[]{microtype}
  \UseMicrotypeSet[protrusion]{basicmath} % disable protrusion for tt fonts
}{}
\makeatletter
\@ifundefined{KOMAClassName}{% if non-KOMA class
  \IfFileExists{parskip.sty}{%
    \usepackage{parskip}
  }{% else
    \setlength{\parindent}{0pt}
    \setlength{\parskip}{6pt plus 2pt minus 1pt}}
}{% if KOMA class
  \KOMAoptions{parskip=half}}
\makeatother
\setlength{\emergencystretch}{3em} % prevent overfull lines
\providecommand{\tightlist}{%
  \setlength{\itemsep}{0pt}\setlength{\parskip}{0pt}}
\usepackage{bookmark}
\IfFileExists{xurl.sty}{\usepackage{xurl}}{} % add URL line breaks if available
\urlstyle{same}
\hypersetup{
  hidelinks,
  pdfcreator={LaTeX via pandoc}}

\author{}
\date{}

\begin{document}

\section{Response to Reviewer 2}\label{response-to-reviewer-2}

\textbf{Manuscript:} \emph{Proxy-Based Diagnostics of Quantum Oracle
Sketching Robustness for Non-IID Sensor and Telemetry Streams}
(submitted as \emph{Quantum Oracle Sketching for Non-IID Sensor and
Telemetry Streams: A Proxy-Based Computational Study})\\
\textbf{Journal:} Sensors (MDPI) · \textbf{Manuscript ID:}
sensors-4470240

We thank both reviewers for careful, constructive reviews that have
materially improved the paper. Both reviews converge on the same central
point, and we agree with it: the previous version, although internally
hedged, still presented itself as evidence about quantum performance,
when what it can honestly support is a \textbf{proxy-based diagnostic
and hypothesis-generation framework for correlation-sensitive streaming
analysis}. The revision repositions the manuscript accordingly ---
title, abstract, statistics, figures, and conclusion --- without
removing any experiment or altering any measured result: all stochastic
outputs regenerate bit-identically from the same seeded pipeline (we
verified this explicitly before re-analysis; the verification scripts
ship with the revision).

\textbf{Summary of the principal changes} (full detail under each point
below):

\begin{enumerate}
\def\labelenumi{\arabic{enumi}.}
\tightlist
\item
  \textbf{Retitled} to \emph{Proxy-Based Diagnostics of Quantum Oracle
  Sketching Robustness for Non-IID Sensor and Telemetry Streams}; the
  abstract now opens ``In theory, \ldots{}'', states that no quantum
  algorithm is implemented or simulated, and was shortened to
  \textasciitilde220 words. \emph{(Note on the tracked-changes file: the
  MDPI class places the title and abstract in the document preamble,
  which latexdiff does not mark; both are completely rewritten.)}
\item
  \textbf{New Section 3.3, ``Relationship Between Operational Proxies
  and Quantum Oracle-Sketching Theory''}, separating established theory
  / assumptions / heuristic constructions / scope limitations, with a
  provenance table (Table 1) stating per quantity what is inherited,
  what is introduced, and what is \emph{not} theoretically guaranteed
  --- including the explicit statement that under the framework's own
  r-only accounting no proxy--baseline crossover exists, so the
  τ-penalised landscape is a classical effective-sample stress test, not
  the framework's cost model.
\item
  \textbf{Statistics rebuilt around estimation rather than testing}
  (Experiment 5): Hodges--Lehmann paired-difference estimates with exact
  signed-rank 95\% intervals, matched-pairs rank-biserial correlations
  (interpreted as sign-concordance), p-values demoted to descriptive
  companions, disclosure that no comparison survives Holm correction at
  n = 10 (with the reason), the τ-estimator lattice sensitivity, and an
  explicit statement that the crossover location is a deterministic
  solution whose uncertainty comes from the τ estimator and the
  constants, not from tests.
\item
  \textbf{New sensitivity analyses from existing data} (no new
  stochastic runs): a C--δ grid for the crossover (Table 8: rho* moves
  from 0.38 to beyond the sweep; ordinal reading is the robust one) and
  an encoder/resolution sensitivity table for the real streams (Table:
  10 re-encodings; the NYC-above-Machine-Temperature ordering holds in
  all 10).
\item
  \textbf{New Discussion subsections} on why analytic and empirical
  correlation estimates diverge, and on what the landscape can and
  cannot predict --- including three \textbf{falsifiable hypotheses
  (H1--H3) with explicit refutation criteria}.
\item
  \textbf{New ``Threats to Validity'' section} (construct / internal /
  external / statistical-conclusion) and a promoted, expanded
  \textbf{Future Work} section.
\item
  \textbf{Conclusion rewritten}: primary contribution = the (τ, r)
  landscape; secondary = proxy-based hypothesis generation; the sentence
  ``The paper demonstrates no quantum advantage and implements no
  quantum algorithm'' appears verbatim.
\item
  \textbf{Figures}: the Experiment-5 legend entry ``Proxy-model
  advantage'' renamed ``Proxy-estimated separation''; in-image labels
  ``n = 10 qubits'' / ``Quantum O(n)'' removed from Figures 1, 2, 6;
  captions aligned with the proxy framing; Figure 8(d)'s synthetic
  markers disclosed as schematic; typographic and overfull fixes
  throughout.
\item
  \textbf{Reference audit (preprints and metadata)}: all 54 references
  were classified, and every entry was then verified at the metadata
  level --- all 51 DOIs were resolved through doi.org content
  negotiation and machine-compared against the cited title, first
  author, venue, and year, with the non-DOI entries checked at source.
  Two corrections resulted: (i) Kallaugher, Parekh and Voronova (SOSA
  2025) now carries the SIAM proceedings DOI (10.1137/1.9781611978315.2)
  instead of its arXiv link, with the printed byline order confirmed on
  the publisher's pages; and (ii) a transposed article number in the Su,
  Guo and Wang (2024) entry was corrected (Information Sciences 676,
  article 120799, DOI 10.1016/j.ins.2024.120799 --- the previously cited
  120783 resolves to an unrelated article). One preprint is retained and
  explicitly flagged: Zhao et al.~(arXiv:2604.07639) is the framework
  this paper studies; we verified via the arXiv record, the authors'
  official listing, and a Crossref registry query that no peer-reviewed
  version exists to date, the entry is labelled ``Preprint'', and every
  invocation of its theorems is co-cited with peer-reviewed anchors
  (Kallaugher, FOCS 2021; Huang et al., Science 2022; Kallaugher et al.,
  SOSA 2025). A citation--bibliography consistency check confirms all 54
  in-text keys map one-to-one onto the 54 entries. Following the
  Academic Editor's request that preprints not be cited, that single
  remaining entry now also carries an explicit footnote at its first
  citation in the Introduction, stating its preprint status and why no
  conclusion of this paper depends on its final published form; a
  separate response to the Academic Editor accompanies this letter.
\end{enumerate}

A tracked-changes PDF (\texttt{tracked\_changes.pdf}, latexdiff against
the reviewed version) accompanies this letter.

\begin{center}\rule{0.5\linewidth}{0.5pt}\end{center}

\subsection{Point-by-point response}\label{point-by-point-response}

We thank the reviewer for an unusually thorough review; the closing
recommendation --- make the (τ, r) landscape-as-hypothesis-generator the
paper's main message --- is precisely the repositioning this revision
performs.

\textbf{R2.1 --- ``The key issue is the relationship between τ and r and
the quantities in the framework of quantum oracle-sketching. \ldots{}
They have to strengthen the theoretical basis for the chosen proxy
construction and distinguish the results from the original from the
assumptions that are used in this paper.''}

\emph{Response.} New §3.3 does this separation explicitly, in four
labelled parts: \emph{Established theory} (the ×R overhead theorem,
O(n)-qubit memory, task-specific O(N\^{}c) hardness --- all results of
the original framework, taken as given); \emph{Assumptions} (two-scalar
summarisability, binarised-stream fidelity,
task-rate-as-pipeline-stand-in --- all modelling choices of this paper,
none inherited); \emph{Heuristic constructions} (the τ division absent
from the framework's accounting; the VC-type rate in place of
oracle-construction cost); and \emph{Scope limitations}. Table 1 gives
per-quantity provenance with an explicit ``not guaranteed'' column. We
also now state the sharpest consequence: under the framework's own
r-only accounting, the Markov, seasonal, and long-memory regimes (r = 1)
incur zero overhead and no proxy--baseline crossover exists --- the
τ-penalised landscape is a classical effective-sample-size stress test,
strictly more pessimistic than the framework's accounting, not a
rendering of it. \emph{Changes:} New §3.3 + Table 1; §3.2;
effective-sample-size lineage now cited (Geyer 1992).

\textbf{R2.2 --- ``The terminology of quantum performance should be
moderated throughout the paper. \ldots{} we will refer to them as
proxy-based predictions instead of the actual quantum performance. This
is particularly relevant in the Abstract, Results, Discussion, figure
captions and Conclusion.''}

\emph{Response.} A systematic terminology pass was made over exactly
those surfaces. ``Quantum-favourable'' no longer appears (now
``proxy-favourable''/``hypothesized favourable''); ``quantum performance
proxy crosses below classical performance'' and similar constructions
were rewritten (``the accuracy proxy computed with the analytic τ lies
below the measured classical baseline''); Table 6's column is ``Proxy
hypothesis''; and --- because captions alone cannot fix text baked into
images --- Figures 1, 2, and 6 were regenerated to remove the in-image
labels ``n = 10 qubits'', ``Quantum proxy (IID, 10 qubits)'', and
``Quantum O(n)''. Explicit ``no quantum algorithm is implemented''
statements now appear in the Abstract, §3.3, §6.5, §6.6, §6.8, §7.4, §8,
and the Conclusion. \emph{Changes:} Global; Figures 1, 2, 6 regenerated;
all captions audited.

\textbf{R2.3 --- ``The statistical comparison between classical
baselines and the proxy model is still not resolved. Our Wilcoxon tests
compare stochastic classical results for seeds to a proxy value that is
deterministic for a given value of parameters. \ldots{} We believe the
magnitude of the classical minus proxy difference and the uncertainty
may be a better gauge than simply p-values.''}

\emph{Response.} Adopted in full --- magnitude and uncertainty now lead,
and the tests are re-scoped to the one thing they can support. Table 4
reports, per sweep point, the Hodges--Lehmann estimate of the paired
difference with its exact signed-rank 95\% interval (e.g., +0.042
{[}+0.011, +0.087{]} at ρ = 0; −0.046 {[}−0.079, −0.003{]} at ρ = 0.88)
and the rank-biserial correlation, which we interpret explicitly as
sign-concordance across seeds (scale-free against a near-constant
reference; all magnitude information lives in the difference estimate).
The p-values are descriptive companions only; we disclose that none
survives Holm correction at n = 10 and why; we disclose the integer-τ
lattice sensitivity (one lag step = 0.012, enough to move the ρ = 0.48
interval across zero); and the interpretation paragraph states that the
crossover is the deterministic solution of acc\_proxy(τ̂(ρ)) = mean
classical accuracy, with uncertainty inherited from the estimator and
the constants (C, δ), not from the tests. The full re-analysis (which
first reproduced the published per-seed results bit-exactly --- all five
p-values to six digits) ships with the revision. \emph{Changes:} Table 4
+ caption; §6.5 (rewritten body and interpretation paragraphs);
Abstract; §7.1; §8; \texttt{verification\_scripts/exp5\_reanalysis.py} +
output.

\textbf{R2.4 --- ``The huge discrepancy between the analytic and
empirical refreshing time is one of the most important findings in the
paper, especially for the long-range dependence regime. \ldots{} The
authors should discuss this in more detail and discuss its implications
for the interpretation of the central (τ,r) sensitivity landscape.''}

\emph{Response.} We agree it is one of the paper's most important
findings and have promoted it accordingly. New §7.3 gives the mechanism
(the analytic proxies target model-level worst-case scales ---
joint-chain mixing; tail-dominated integrated autocorrelation --- while
the empirical estimator measures the short-lag marginal decorrelation of
the binarised stream; for a non-mixing process there is no scalar the
two could agree on), and the implications for the landscape are now
carried through the whole paper: verdicts are bound to their estimator
convention everywhere, and we state openly that the burst and
long-memory extremes \emph{exchange places} between the analytic and
empirical conventions (burst saturates through its measured duplication;
long-memory's high empirical value is the short-lag artefact) --- in
Experiment 9, §7.3, the RQ2/RQ3 answers, and the Figure 3/7 captions.
The abstract now says the two estimates ``diverge'' (not
``overestimate'', which presupposed the empirical value as ground truth)
and adds that they answer different questions. \emph{Changes:} New §7.3;
§6.9; §6.3; Figures 3, 7 captions; Conclusion RQ1--RQ3; Abstract.

\textbf{R2.5 --- ``The classical-memory comparison needs to be taken
with care. The classical lower bound curve \ldots{} is not actually an
empirical lower bound for the SGD, averaged-SGD, or Count-Min tools used
in our experiments. \ldots{} we should make it more prominent in the
dimension-scaling discussion.''}

\emph{Response.} Made prominent exactly there. Experiment 6 now opens
with: ``No memory lower bound is measured empirically in this
experiment: the curves in panel (a) plot cited theoretical quantities,
and no result here constitutes an empirical memory floor for online SGD,
averaged SGD, or Count-Min.'' Figure 6's caption labels panel (a)
``cited bounds; no measured content'' and the in-image curve labels now
read ``(cited theory)''. §3.3's ``Memory references'' paragraph and the
provenance table carry the same statement structurally, and §7.1 states
that none of its structural factors is a measured quantum result.
\emph{Changes:} §6.6; Figure 6 (caption and regenerated labels); §3.3;
Table 1; §7.1.

\textbf{R2.6 --- ``The real-stream analysis is useful as a sanity check
but it is limited to two NAB datasets. The conclusions on the sensor and
telemetry applications should therefore be conservative.''}

\emph{Response.} The conclusions are now explicitly conservative: §6.8
states the experiment ``is illustrative rather than statistically
representative'', the External-validity paragraph records that
generalisation is unverified, and the abstract claims placement only.
Additionally, a new encoder/resolution sensitivity table strengthens
what \emph{can} be said from two streams: across ten re-encodings of the
same raw series (three sampling resolutions × three bucket widths), the
NYC-above-Machine-Temperature ordering holds in every configuration, so
the reported placement is at least encoder-stable. Broader corpora are
committed in Future Work with a concrete test protocol (H1).
\emph{Changes:} §6.8; new encoder-sensitivity paragraph + table; §8; §9;
§7.4 H1.

\textbf{R2.7 --- ``the target defined based on a percentile of the full
time series should also be reconsidered in the context of streaming
analysis, since using the full series to set the target threshold can
take into account future observations. A threshold set from a first
calibration interval or updated only from previously observed samples
would provide a more rigorous prequential design.''}

\emph{Response.} This is an excellent catch, and we now disclose it
explicitly rather than leave it implicit. §6.8 (``Binarisation choice'')
states that the threshold is computed from the full series and therefore
uses information not available prequentially at each step; that the
feature stream and the (τ, r) estimates involve no such look-ahead ---
so the landscape placement, the paper's central output for these
streams, is unaffected, the leakage touching only the classical-baseline
reference numbers; and that a threshold calibrated on an initial
interval, or updated from previously observed samples only, is the
rigorous prequential design, which we adopt as future work. The Threats
section (Internal validity, item (v)) records it as a residual risk. We
chose disclosure over re-running the baselines in this revision to keep
every published number bit-reproducible; we are happy to re-run with a
calibration-interval threshold in a further round if the reviewer
prefers. \emph{Changes:} §6.8 Binarisation paragraph; §8 Internal
validity (v); §9.

\textbf{R2.8 --- ``The interpretation of the NYC Taxi Count-Min result
also needs to be qualified. The six-bit representation gives a
relatively small number of possible feature vectors and hence memorising
the same feature patterns is relatively easy. \ldots{} I would like to
give more weight to the implications when we think about the apparent
agreement between the Count-Min performance and our proxy prediction.''}

\emph{Response.} The qualification now carries real weight in three
places. The dedicated ``Why Count-Min looks strong'' paragraph is
retained; the interpretation adds, plainly: ``the favourable-side
separation between the two streams rests on Count-Min alone; the
hashed-linear learners do not corroborate it'' --- and we additionally
report that both SGD variants fall \emph{below} the majority baseline on
accuracy and that Averaged-SGD's balanced accuracy straddles chance, so
the apparent proxy--Count-Min agreement is explicitly not read as
validation. §7.4 lists representation-specific effects like this
memorisation route among the things the landscape \emph{cannot}
anticipate, and --- going further in the reviewer's direction --- the
same section now concedes that, read strictly as an accuracy projection,
the proxy is falsified on Machine Temperature (0.500 predicted vs
0.72--0.80 observed), with the rare-class diagnostic reading offered as
interpretation, not confirmation. Table 6 now carries the
ordering-to-ordering intervals (Count-Min F1 = 0.644 ± 0.062,
per-ordering range 0.43--0.72), which tempers the headline number.
\emph{Changes:} §6.8 (findings bullets, Count-Min paragraph,
interpretation (ii)--(iii)); Table 6 (intervals restored); §7.4.

\textbf{R2.9 --- ``In particular, the target function ablation should be
highlighted because it shows that our classical performance can change
dramatically with the classification task in question and that our proxy
is in fact target independent. A correlation-based proxy that is
insensitive to the target function cannot necessarily predict the
performance of algorithms for different kinds of decision problems.''}

\emph{Response.} Highlighted and stated as a scope boundary in the
reviewer's own terms. §6.9 (Target function) now reads: ``classical
learner performance is strongly target-dependent, while the proxy
depends only on (τ, r, T) and is target-independent by construction. The
proxy therefore cannot universally predict downstream classifier
performance; it characterises the correlation structure of the stream
--- the effective sample budget available to any bounded-memory learner
--- not task-specific learnability.'' §7.4 elevates this to the
``cannot'' list (with the further empirical point that measured
classical accuracy is flat while T\_eff falls fivefold, so the budget
model does not predict even the classical arm's accuracy), the Figure 9
caption states the proxy ``does not track task-specific learnability'',
and RQ2's answer now says τ and r are budget summaries, not accuracy
predictors. The ablation's protocol (prequential, lr = 0.01, T = 4096)
is also now stated so its absolute levels are not misread against the
main experiments. \emph{Changes:} §6.9; §7.4; Figure 9 caption;
Conclusion RQ2.

\textbf{R2.10 --- ``the strongest contribution of the paper is the (τ,r)
sensitivity landscape proposed as a hypothesis-generating tool \ldots{}
The paper would be better off having this methodological contribution as
the main message rather than the separation of quantum-classical
performance.''}

\emph{Response.} This is now the paper's stated identity, end to end:
the title leads with ``Proxy-Based Diagnostics''; the abstract announces
``a proxy-based diagnostic and hypothesis-generation framework''; the
contribution list names the landscape ``the paper's central diagnostic
and hypothesis-generation artifact''; §7.4 states the falsifiable
hypotheses it generates; and the Conclusion opens by declaring the
landscape the primary contribution and hypothesis generation the
secondary one, followed by: ``The paper demonstrates no quantum
advantage and implements no quantum algorithm.'' \emph{Changes:} Title;
Abstract; §1; §7.4; §10.

\begin{center}\rule{0.5\linewidth}{0.5pt}\end{center}

\subsection{Note on verification}\label{note-on-verification}

Before making any statistical change we reproduced the published
Experiment-5 per-seed results bit-exactly from the pinned seeds (all
five Wilcoxon p-values match the published values to six significant
digits), and all figure regenerations are rendering-only (numerical
outputs byte-identical). The re-analysis and sensitivity scripts, with
their outputs, are included in the accompanying package
(\texttt{verification\_scripts/exp5\_reanalysis.py} and
\texttt{verification\_scripts/encoder\_sensitivity.py};
\texttt{code/exp5\_effectsizes.py} in the public repository), and the
reproducibility documentation was corrected and extended (exact seed
indices per experiment; the previously missing Machine-Temperature
download step).

On behalf of all authors,

\textbf{AbdelMoniem Helmy} (corresponding author)
abdelmoniem.hafez@cu.edu.eg · ORCID 0000-0001-5996-6019

\end{document}
